% Corte mecánico íntegro: parte_ii.tex:323--419.
% Residencia material del ensamblaje sucesor de 102 capítulos.
% Residencia causal: capítulo 46.

\chapter{Vis spectrale de Pell et suspension logarithmique}

\section{Loi de vis}

Le corpus contient l'identité exacte
\begin{equation}
 V^{-1}\mathscr S V=(2+\sqrt3)\ii\,\mathscr S.
 \label{espiral:eq:tornillo-pell}
\end{equation}
Soit \(\varepsilon=2+\sqrt3\). L'orbite transverse
\[
 z_n=z_0(\varepsilon\ii)^n
\]
satisfait
\[
 r_n=r_0\varepsilon^n,
 \qquad
 \theta_n=\theta_0+\frac{\pi n}{2}.
\]
Eliminando \(n\),
\begin{equation}
 r(\theta)=r_0
 \exp\!\left[
 \frac{2\log(2+\sqrt3)}{\pi}
 (\theta-\theta_0)
 \right].
 \label{espiral:eq:espiral-log-pell}
\end{equation}
La spirale logarithmique est une sortie de l'opérateur : son équation n'a pas été choisie pour ajuster ensuite \(\varepsilon\).

\begin{proof}[Dérivation de \cref{espiral:eq:espiral-log-pell}]
De \(\theta_n-\theta_0=n\pi/2\) découle \(n=2(\theta_n-\theta_0)/\pi\). En substituant dans \(r_n=r_0\varepsilon^n\),
\[
 r_n=r_0\exp\left(
 \frac{2\log\varepsilon}{\pi}(\theta_n-\theta_0)
 \right).
\]
L'interpolation continue est l'équation indiquée.
\end{proof}

\section{Involution conjuguée de contraction}

L'extension bilatérale \(n\in\Z\) produit
\[
 n\to+\infty:\ r_n\to\infty,
 \qquad
 n\to-\infty:\ r_n\to0.
\]
L'inversion \(n\mapsto-n\) échange expansion et contraction. Comme \(\varepsilon^{-1}=2-\sqrt3\), les deux branches appartiennent à la même unité de Pell. Ce ne sont pas deux spirales indépendantes, mais deux orientations de la même vis spectrale.

\begin{figure}[htbp]
\centering
\begin{tikzpicture}[scale=.82]
 \draw[->,HMTGray] (-5.1,0)--(5.1,0) node[right] {$x$};
 \draw[->,HMTGray] (0,-3.1)--(0,3.1) node[above] {$y$};
 \draw[HMTBlue,very thick,domain=-7:5.2,samples=240,smooth,variable=\t]
   plot ({.27*exp(.22*\t)*cos(deg(\t))},{.27*exp(.22*\t)*sin(deg(\t))});
 \draw[HMTRed,thick,domain=-7:5.2,samples=240,smooth,variable=\t]
   plot ({.27*exp(.22*\t)*cos(deg(-\t))},{.27*exp(.22*\t)*sin(deg(-\t))});
 \foreach \k in {-4,-3,...,4}{
   \pgfmathsetmacro{\ang}{1.57079632679*\k}
   \fill[HMTNavy]
     ({.27*exp(.22*\ang)*cos(deg(\ang))},
      {.27*exp(.22*\ang)*sin(deg(\ang))}) circle (.8pt);
 }
 \node[HMTBlue,font=\small\sffamily] at (3.5,2.4) {expansion};
 \node[HMTRed,font=\small\sffamily] at (3.5,-2.4) {orientation conjuguée};
\end{tikzpicture}
\caption{Projection transverse de la vis de Pell. Les marques représentent des quarts de tour.}
\label{espiral:fig:tornillo-pell}
\end{figure}

\section{Covariance d'échelle et conservation énergétique}

La dilatation \(\varepsilon\) ne doit pas être interprétée comme une croissance gratuite de l'amplitude énergétique d'une onde dans le vide. Sa réalisation cohérente agit sur l'échelle spatiale ou spectrale :
\[
 (r,\varphi,\lambda)
 \longmapsto
 (\varepsilon r,\varphi+\pi/2,\varepsilon\lambda).
\]
Cette distinction sera essentielle dans la réalisation électromagnétique de la \cref{espiral:chap:realizacion-em}.

\begin{resultbox}[title={Compatibilité exacte entre la vis spectrale de Pell et la suspension phase--mémoire}]
Le relèvement phase--mémoire minimal est une hélice circulaire. La vis spectrale combine un quart de tour et une dilatation de Pell ; sa projection est une spirale logarithmique. Tous deux proviennent du registre HMT, mais ne sont pas la même périodicité et ne s'obtiennent pas l'un à partir de l'autre par oubli des types.
\end{resultbox}
